% UBT-AI-PROVENANCE-BEGIN
% schema: ubt-ai-provenance/v1
% tier: C_working
% ai_assistance: disclosed
% human_review: risk-based
% editorial_responsibility: Ing. David Jaroš
% policy: ../../../AI_PROVENANCE.md
% notice: Working material; exhaustive human review is not claimed.
% UBT-AI-PROVENANCE-END

\chapter{Tests, reproducibility and formal verification}
\label{ch:tests-repro}

The live verification layer is the active repository, not \texttt{ARCHIVE/}.
Important entry points are:
\begin{itemize}
 \item \texttt{tests/} for scientific, provenance and architectural regression tests;
 \item \texttt{tools/apply\_provenance\_headers.py} for AI provenance marking;
 \item \texttt{tools/latex\_audit.py} for isolated translations of LaTeX documents;
 \item \texttt{tools/check\_alpha\_claims.py} and active alpha tests for
       review of claims sensitive to current status;
 \item numerical checks of theta/Mellin/Gaussian branches are provided by\\
       \texttt{research\_tracks/theta\_spectral/}\\
       \texttt{theta\_zeta\_probe.py};
 \item \texttt{docs/DERIVATION\_VERIFICATION\_POLICY.md} as binding rules
       derivation checks in papers.
\end{itemize}

\section{Four different types of control}
It is important not to confuse the different levels of verification. LaTeX translation only proves
that the document is syntactically composable. The numerical test verifies the specific
numerical consequences in a given accuracy. A CAS tool, such as SymPy or Maxima,
it verifies the token identity we actually gave it. Formal system
Lean checks whether a formalized conclusion follows from formalized assumptions.
Neither of these checks alone proves that the relevant assumptions are
physically selected by canonical UBT dynamics.

\section{Lean as the preferred formal controller}
It is the preferred tool for new or substantially changed theorem derivations
Lean. If the statement is reasonably formalizable, the goal is to have alongside the analytic
derivation also compiling \texttt{.lean} proof. If there is no formalization yet
done or is not available in a given Lean environment, the status must be explicit
marked as \texttt{LEAN-PENDING}; one must not write that the statement is ``formally
proven'' just because the agent generated the Lean syntax.

Lean does not replace other controls. For the main result of the paper, the second is desirable,
independent pipeline: for example SymPy/Maxima for symbolic algebra and
NumPy/Octave/MATLAB/Mathcad for numerical reproduction. Two implementations should
if possible use different wording so as not to repeat the same mistake.

\section{Mandatory record of paper verification}
Each new or substantially changed paper should contain a section\emph{Verification}
or a link to an accompanying record. For the main statement, at least:
\begin{itemize}
 \item statement, equation, or lemma identifier;
 \item analytical status and assumptions;
 \item used tool and its version;
 \item path to reproduction script, test, worksheet or \texttt{.lean}
       file;
 \item result and possible tolerance;
 \item exactly what the check verifies and what it does not;
 \item Lean formalization status: \texttt{PROVED}, \texttt{PARTIAL},
       \texttt{LEAN-PENDING}, or \texttt{NOT-APPLICABLE}.
\end{itemize}

\section{CI textbook}
The dedicated workflow is \texttt{.github/workflows/textbook.yml}. It runs on
relevant pushes and pull requests and can also be started manually. It first
checks bilingual parity and then uses \texttt{tools/latex\_audit.py --strict}.
A failure in any standalone textbook document therefore fails the workflow.

After the strict audit, the workflow calls the same Makefile used locally and
creates exactly two official PDFs in \texttt{build/textbook/public/}:
\begin{itemize}
 \item \texttt{UBT\_Textbook\_EN.pdf} --- English edition;
 \item \texttt{UBT\_Textbook\_CS.pdf} --- Czech edition.
\end{itemize}
Each PDF is uploaded as a separately named GitHub Actions artifact. Email
delivery attaches both editions and is allowed only after a successful push
to \texttt{master}; pull requests and manual runs never send email.

For local published copy can be run
\begin{verbatim}
make -C docs/textbook publish
\end{verbatim}
which copies the same pair into \texttt{docs/pdfs/}. Generated PDFs are only
publication outputs; scholarly status remains governed by source text and the
provenance/status registry.

\section{Minimal local check before textbook patch}
The recommended sequence is:
\begin{verbatim}
python tools/apply_provenance_headers.py --apply
python tools/apply_provenance_headers.py --check
python tools/latex_audit.py --include 'docs/textbook/*.tex' \
  --include 'docs/textbook/**/*.tex' --jobs 2 --timeout 240 --strict
make -C docs/textbook verify
\end{verbatim}
A successful build confirms syntactic and reproduction properties. It does not increase itself from
self-Tier-C claims on proven UBT result.
